\resheading{Research Highlights}
     \ressubheading{}{}{ \bf High-performance computing and artificial intelligence at ANL} {Jan 2014 - Present}
\begin{itemize}

\resitem{Application and evaluation of large language models (LLMs) for scientific tasks.}
\begin{itemize}
\resitem{Developed an agentic framework for computational chemistry simulations.}
\resitem{Created an LLM evaluation benchmark for computational chemistry and initiated a HuggingFace leaderboard.}
\resitem{Implemented a Retrieval-Augmented Generation (RAG) pipeline for ALCF documentation and evaluated different large-language models.}
\resitem{Generated (both manually and automated) multiple-choice questions for chemistry, physics, and material science domains for LLMs.}
\resitem{Evaluated different retrieval augmentation and fine-tuning methods to customize LLMs for domain-specific tasks, enhancing their accuracy and relevance in generating contextually appropriate responses.}
\resitem{Explored utilization of LLM agents for enhancing scientific workflows, focusing on their ability to automate data analysis, streamline complex computational tasks, and improve the overall efficiency. }
\end{itemize}
\resitem{Argonne point-of-contact and a co-developer for the \href{https://github.com/NWChemEx}{NWChemEx Project}, an ECP funded application software for quantum chemistry. }
\begin{itemize}
    \resitem{Contributed to the implementation of the C++ infrastructure code responsible for caching and serialization mechanisms.}
    \resitem{Contributed to the design and implementation of the user-interface based on Python bindings.}
    \resitem{Actively participating in the development and maintenance activities such as adding features, fixing issues, reviewing pull requests, and improving documentation.}
\end{itemize}

\resitem{Developing machine learning potentials for material science and chemistry.}
\begin{itemize}
    \resitem{Developed Gaussian Approximation Potentials (GAPs) for two-dimensional materials and verified their accuracy for thermal properties. Published the models and training data on \href{https://github.com/AI4Materials/2DM-GAP}{GitHub}.}
    \resitem{Compared the accuracy and training efficiency of GAPs with neural network based potentials.}
    \resitem{Implemented distributed training for equivariant graph neural networks, benchmarked and optimized the performance on Polaris.}
    \resitem{Developed neural network potentials for salt mixtures using density functional theory calculations.}
\end{itemize}
\resitem{Co-developer of \href{https://github.com/VALENCE-software/VALENCE}{VALENCE} quantum chemistry code based on variational subspace valence bond method.}
\begin{itemize}
    \resitem{Developed Vtools that can automate input preparation and visualization.}
    \resitem{Contributed to optimization and benchmarking of the code on Intel-KNL nodes of Theta supercomputer.}
\end{itemize}
\resitem{Reduced training time of image segmentation code for connectomics.}
%\resitem{Member of ALCF data science group on research, development and support in machine learning to glean insights from large-scale datasets being produced by computational simulations, experiments and observations on high performance computing systems.}
\begin{itemize}
\resitem{Enabled and optimized data parallel training of deep convolutional neural networks for 3D image segmentation by using Horovod and hyperparameter optimization.}
\resitem{Reduced the training time from seven days to less than 5 hours by using 1024 Intel KNL CPUs on ALCF Theta supercomputer.}
\resitem{Parallelized the preprocessing step with mpi4py to reduce the computational time from 4 hours to five minutes.}
%\resitem{Developed a notebook to use Balsam workflow tool on ALCF Jupyter Hub to submit and monitor jobs using only a browser.}
%\resitem{Integrated sparse eigensolvers (slepc4py) for a manifold learning package to improve computational performance. See \url{https://github.com/keceli/megaman}.}
\end{itemize}
%\resitem{Development and implementation of a tight-binding DFT and SEMO code based on a sparse direct eigensolver. Analyzed the strong scaling of this code beyond 200,000 BG/Q cores for problems involving more than 100,000 atoms. The code is built on PETSc and SLEPc libraries.}
\resitem{Co-developer of \href{https://github.com/Auto-Mech/}{Auto-Mech} software suite for computational thermochemistry and kinetics.}
\begin{itemize}
\resitem{Developed \href{http://github.com/keceli/qtc}{Quantum Thermochemistry Calculator}} to automate complex workflows for thermochemistry calculations. 
\resitem{Working on the curation and optimization of a database to efficiently store and manage the results generated by Auto-Mech.}
\resitem{Working on porting Auto-Mech codes to Polaris and Aurora and optimizing MESS code with multiprecision algorithms.}
\end{itemize}
\resitem{Development, testing and benchmarking of sparse and dense eigensolvers for quantum chemistry applications.}
\begin{itemize}
    \resitem {Developed shift-and-invert spectrum slicing eigensolver (SIPs) based on PETSc and SLEPc libraries.}
    \resitem {Demonstrated strong scaling beyond 200,000 cores and enabled density functional tight-binding (DFTB) simulations of more than 100,000 atoms. This work was published on the cover of Journal of Computational Chemistry and highlighted on ALCF newsbytes.} 
    \resitem{Integrated the massively parallel sparse eigensolver, SIPs with \textsc{SIESTA} ab initio molecular dynamics package and demonstrated higher strong scaling efficiency and reduced time to solution compared to default solver, ScaLAPACK.The code has been ported and tested on BG/Q nodes at ALCF-Mira, Intel-Haswell nodes at NERSC-Cori, and Intel-KNL nodes at ALCF-Theta supercomputers.}
    \resitem{Developed a Python code, PSCF, based on mpi4py, petsc4py and slepc4py for semi-empirical molecular orbital calculations.}
\end{itemize}
\resitem{Performed benchmark coupled-cluster and DFT calculations with NWChem and Molpro for accurate thermochemistry of transition metal oxide clusters. Some of these calculations involve more than 1300 basis functions and were run on supercomputers.}
\end{itemize}

\ressubheading{}{}{\bf Chemical engineering and \bf software development at MIT} {July 2012 - Jan 2014}
\begin{itemize}
\resitem{Developed a ``fit-for-purpose" model using kinetic model reduction and  lumping analysis in a proprietary hydrotreating reactor model as part of BP-MIT Advanced Conversion Research Program. The reductions in the model can eliminate various feed characterization experiments and reduce capital costs significantly.}
%\resitem{Developed a two-phase (liquid and vapor) plug-flow hydrotreater reactor model to test different lumping schemes.}
%\resitem{Identified the mathematical connections of kinetic mechanism reduction methods such as affine lumping, partial equilibrium assumption, quasi steady state approximation, method of invariant manifold and represented them using a projection-based method.}
\resitem{Served as a member of the software development team of reaction mechanism generator (RMG) program, which constructs kinetic models composed of elementary chemical reaction steps. Contributed to the coding, testing and bug-fixing efforts before the major release of the Java version and lead the parallelization of the Python version. See \url{http://github.com/ReactionMechanismGenerator/RMG-Java}. }
%Code repository is  at \url{https://github.com/GreenGroup}.}
%\begin{itemize}
 %  \resitem{Contributed to the coding, testing and bug-fixing efforts before the major release of the Java version and leading the parallelization of the Python version using the SCOOP module.}
%   \resitem{William H. Green, Joshua W. Allen, Beat A. Buesser, Robert W. Ashcraft, Gregory J. Beran, Caleb A. Class, Connie Gao, C. Franklin Goldsmith, Michael R. Harper, Amrit Jalan, {\bf{ M. Ke\c{c}eli}}, Gregory R. Magoon, David M. Matheu, Shamel S. Merchant, Jeffrey D. Mo, Sarah Petway, Sumathy Raman, Sandeep Sharma, Jing Song, Yury Suleymanov, Kevin M. Van Geem, John Wen, Richard H. West, Andrew Wong, Hsi-Wu Wong, Paul E. Yelvington, Nathan Yee, Joanna Yu; “RMG - Reaction Mechanism Generator v4.0.1”, 2013.}
%   \end{itemize}
\end{itemize}
